%% kapa: exam=1 round=35 year=2024 period=1 form=A subject=economics subject_ko=경제학원론 seq=24 number=64
%% source: Q-Net 감정평가사 기출문제 (cst00309) artlSeq=5213229 file=raw/1008/5213229/1교시 문제지 원본.pdf
%% conversion: kordoc 4.24.51 -> tools/gen1.py (auto); review=proofread-sonnet; flags=none
\begin{question}{64}
필립스곡선이 단기에는 우하향하고 장기에는 수직인 경제에서 중앙은행은 테일러 준칙(Taylor's rule)에 의해 통화정책을 시행한다. 중앙은행이 높은 인플레이션율을 낮추기 위해 인플레이션 감축정책(디스인플레이션 정책)을 시행할 때, 이에 관한 설명으로 옳은 것을 모두 고른 것은?
\begin{framedtext}
ㄱ. 기대인플레이션이 빨리 조정될수록 장기균형에 빨리 도달한다. \par
ㄴ. 단기에는 실질이자율이 하락한다. \par
ㄷ. 단기에는 총생산이 감소하여 경기침체가 나타난다. \par
\end{framedtext}
\begin{choices}
  \item ㄱ
  \item ㄴ
  \item ㄱ, ㄷ
  \item ㄴ, ㄷ
  \item ㄱ, ㄴ, ㄷ
\end{choices}
\end{question}
\answer{3}
